% Paper sections as of 8 Nov 2026 (copied from the NAMRC 55 manuscript).
% Requires booktabs, graphicx and the bibliography entries bednarik2018tds, ronneberger2015unet, shi2015convlstm.
\section{Methodology}
This study developed and evaluated the perception component of a framework for printing on a deforming surface: estimating the surface where it is hidden by the print head and forecasting its height over a short horizon. Material was not deposited. Instead, a servo-actuated substrate generated programmable surface motion while a depth camera observed it, and a print head placed in the camera view reproduced the occlusion that occurs during deposition. The work comprised the experimental setup and depth acquisition, the construction of region-of-interest height maps, and the design, training, and evaluation of an occlusion-aware surface reconstruction and forecasting model.

\subsection{Theoretical Background}
The receiving surface was represented as a time-dependent height field,
\begin{equation}
    z = S(x,y,t),
\end{equation}
where $S(x,y,t)$ denotes the surface height at planar location $(x,y)$ and time $t$. This representation assumes that each planar
location corresponds to a single visible surface height. The formulation is therefore restricted to predominantly vertical
deformation, gentle local slopes, and negligible lateral material motion. Overhangs and substantial lateral displacement were outside the initial formulation. Regions hidden from the camera by the print head were addressed by the occlusion-aware model described below.

Depth measurements were converted into three-dimensional camera coordinates using the corresponding stream intrinsics.
For an undistorted pinhole observation with depth $d$ at pixel coordinates $(u,v)$,
\begin{equation}
    \mathbf{p}_{c} =
    \begin{bmatrix}
        (u-c_x)d/f_x \\
        (v-c_y)d/f_y \\
        d
    \end{bmatrix},
\end{equation}
where $f_x$ and $f_y$ are the focal lengths and $(c_x,c_y)$ is the principal point. Intrinsics were taken from the factory calibration recorded with each session, and depths were expressed in millimeters.

The resulting camera points were aggregated onto a regular camera-XY grid, with each supported cell storing the mean camera-axis depth, a validity indicator, and its supporting point count. Heights in this study therefore represent camera-axis depth rather than height in a printer coordinate system. Real print-head comparisons were additionally re-expressed as vertical differences in the printer frame using a provisional camera-to-printer registration, described below.

In a printing application, the surface estimate is needed when the nozzle reaches a deposition location rather than when the surface is observed. The required prediction horizon is
\begin{equation}
    \tau = t^{\mathrm{use}}-t^{\mathrm{observation}},
\end{equation}
which includes observation age, processing and communication delays, and tool travel. The model was therefore designed both to reconstruct the current surface and to forecast it over a short horizon.

\subsection{Experimental Setup}
The testbed reproduced the sensing conditions of printing on a moving surface without depositing material. A servo-actuated substrate generated programmable surface motion, and an Intel RealSense D405 depth camera observed the surface from a fixed pose. A Raspberry Pi 5 commanded the servo through 50~Hz pulse-width modulation and recorded the camera streams. To reproduce the occlusion caused by a nozzle during deposition, the recordings used for real-occlusion evaluation placed a stationary print head within the camera's view of the surface. Models were trained and evaluated on a workstation with an NVIDIA GeForce RTX 5090 GPU (PyTorch 2.11 with CUDA 12.8), except for the 30-epoch reference model described below, which was trained earlier on a different computer.


\subsection{Servo-Driven Surface Recordings}
Occlusion-aware reconstruction and forecasting were developed with recordings of a surface deformed by a servo under programmed motion. Twenty motion patterns were recorded with the D405 at a resolution of $640 \times 480$ pixels and a nominal rate of 30 frames per second, for approximately 30~s each. The 20 recordings contained 18,019 paired color and depth frames; delivered color frame rates ranged from 29.93 to 30.03 frames per second, and the 18 camera frame-number gaps in each stream were preserved rather than interpolated. Color and depth were stored losslessly with the FFV1 codec, and every decoded frame was verified against a per-frame SHA-256 checksum. Native depth and color frames were saved without alignment; recorded depth-camera intrinsics, depth scale, factory depth-to-color extrinsics, camera timestamps, and host monotonic timestamps were retained.

The patterns comprised constant-speed, sinusoidal, polynomial, and asymmetric sweeps; sweeps with holds; a staircase; seeded random waypoints; small, medium, expanding, and contracting oscillations; rising and falling chirps; superimposed and damped oscillations; center-out excursions; double pulses; stop--start motion; and an oscillation about a shifting center. Commanded servo angles remained within 30--270$^\circ$, with each pattern's amplitude reduced where necessary to respect a 30$^\circ$/s speed limit, and commands were updated at a nominal interval of 20~ms. The servo labels represented commanded positions at host receipt and were not exposure-synchronized measurements of the actual servo angle.

Five additional patterns---amplitude beats, slow--fast--slow oscillation, three plateaus, a seeded random walk, and an interrupted chirp---were reserved for evaluation with a real print head. Each was recorded twice with the same commanded trajectory and camera settings: once with the print head in the camera view, and once approximately 45 minutes earlier without it. None of these recordings was used for training, validation, or model selection.


\subsection{Region-of-Interest Height Maps}
Native depth values were converted to millimeters using the recorded depth scale, approximately 0.1~mm per unit. Zero depth values were treated as invalid. To select a conservative central portion of the surface, depth pixels were restricted to $245 \leq u < 420$ and $205 \leq v < 315$, and camera-axis depths were restricted to 230--450~mm. The selected depth pixels were unprojected using the depth-stream intrinsics and aggregated by mean depth onto a $23 \times 45$ camera-XY grid with 4~mm cells; each valid cell contained at least one accepted point. The resulting maps describe camera-axis depth, not height in the printer coordinate system. The image crop and depth gate were scene-specific approximations rather than verified surface segmentation and did not include the entire substrate.


\subsection{Model Selection}
The model was required to perform two coupled tasks from the same partially observed history: reconstruct surface cells currently hidden by the print head, and forecast the whole surface over a short horizon. Three properties of the problem guided the choice of architecture. First, the observations are regular height maps in which neighboring cells deform together, so a hidden cell can be inferred in part from the visible surface around it. Second, the surface moves continuously, so a cell that is hidden now may have been observed a few frames earlier, and its recent motion constrains its current and future height. Third, the occluded region is large relative to the grid and changes over time, so neither spatial nor temporal information alone is sufficient.

Simple alternatives address only part of this problem. Carrying forward the last observation, or fitting velocity or acceleration to each cell independently, uses no spatial context and cannot estimate a cell that was not observed within the input window. Per-frame spatial interpolation or inpainting fills hidden regions from the current frame alone and discards the surface's motion. Fully connected recurrent networks treat the grid as an unordered vector and lose the spatial structure of the height map.

Convolutional long short-term memory (ConvLSTM) replaces the matrix products of a conventional LSTM with convolutions, so that its hidden and memory states remain spatial maps~\cite{shi2015convlstm}. Each cell's state is updated from its spatial neighborhood and from its own history, which suits a deforming height field. Its gated memory can retain an estimate for a cell through frames in which that cell is hidden, and its latent state can be advanced without new observations to produce successive forecasts, which provides the multi-step rollout used here.

A ConvLSTM operating at the full grid resolution has a receptive field limited by its kernel size, whereas the masked region can span many cells. A U-Net encoder--decoder~\cite{ronneberger2015unet} addresses this limitation. The encoder downsamples the $23 \times 45$ grid to a $6 \times 12$ latent map, where the recurrent cells operate and where a $3 \times 3$ kernel covers a large fraction of the surface. The decoder restores the full resolution, and its skip connections return fine spatial detail from observed regions that would otherwise be lost through downsampling. The recurrent latent state therefore captures broad, slowly varying deformation, while the skip connections preserve local detail where the surface is visible.

The observation mask, observation age, and never-observed indicator are supplied as explicit input channels, so that the network distinguishes a measured height from a placeholder and can weight older observations accordingly. Because reconstruction and forecasting share one latent state, an estimate of a hidden cell benefits from the temporal context that also drives the forecast. The resulting network contains 708,753 trainable parameters, which is compact relative to the available training data. It also processes frames causally, so it can be applied to a live depth stream without future information.

Two initializations were compared. In the first, the network was trained from randomly initialized parameters using only the servo-driven recordings, so that its parameters were fitted to the specific camera pose, surface, region of interest, and motion range of this setup. In the second, the same network was first pretrained on the rectangular-cloth subset of the EPFL texture-less deformable surfaces dataset~\cite{bednarik2018tds}, which contains 18 Kinect depth sequences of a deforming cloth, and then fine-tuned on the servo recordings. Because the network is fully convolutional and its input channels do not depend on grid size or frame rate, the pretrained weights could be loaded unchanged despite the different grid and cadence of the cloth data. In either case, the trained model is specific to this scene; a different camera pose, substrate, or grid requires new recordings and retraining.


\subsection{Occlusion-Aware Surface Reconstruction and Forecasting}
Complete motion-pattern recordings were assigned to partitions before temporal windows were generated. Sixteen recordings were used for training, two for validation, and two for testing. Validation used the two-frequency and damped-oscillation recordings, and testing used the stop-start and biased-oscillation recordings. The recordings shared the same experimental setup and surface; the test partition therefore evaluated held-out motions in that setup. A window contained 25 input frames, spanning approximately 800~ms, followed by six target frames at nominal 33.3~ms intervals, so that the final target corresponded to a nominal 200~ms forecast horizon. Overlapping training windows used a starting-frame stride of 12, whereas validation and test windows used a stride of one. Windows containing any interval outside a $\pm20\%$ tolerance of 33.3~ms, or any frame with less than 10\% valid grid coverage, were excluded. This yielded 1,128 training windows, 1,744 validation windows, and 1,740 test windows.

Synthetic occlusion was applied only to copies of the input histories. For each window with a standard mask, an anisotropic Gaussian was constructed on the grid with a randomly selected orientation and an aspect ratio between 0.6 and 1.7. The 20\% of grid cells with the highest Gaussian values formed a binary mask. Its center moved by at most one cell along each axis every six input frames (approximately 200~ms), reflecting at grid boundaries, while its orientation and aspect ratio remained fixed within the window. Thus, nearby cells were blocked together, while the Gaussian values themselves did not replace or blur the measured depths. The mask size was specified relative to the entire grid, so its overlap with valid surface cells could vary. The original valid measurements provided the current and future reference targets; naturally invalid measurements were excluded from supervision. Masks were deterministic for a given seed, session path, and window, with separate seeds for training, validation, and testing.

Because a print head close to the camera can hide most of the region of interest, a head-scale mask was also defined. It used the same Gaussian construction, but its masked fraction was drawn uniformly between 40\% and 90\% of the grid for each window, and in half of the windows it was held stationary throughout the input history; otherwise it moved as above. Models were trained either with standard masks alone or with mixed masks, in which each window received a head-scale mask with probability 0.5 and a standard mask otherwise. Validation masks were drawn from the same distribution as the training masks.

The occlusion-aware model combined a U-Net encoder--decoder with convolutional long short-term memory (ConvLSTM) cells. Each input frame contained four channels: normalized observed depth, a binary observation mask, time since the cell was last observed within the input window, and an indicator for cells not previously observed in that window. Observation age was capped and normalized using a 1,000~ms scale; missing values were represented by finite placeholders. Depth normalization used the mean and standard deviation of valid training measurements only. The encoder extracted spatial features at full, half, and quarter resolution, while recurrent dynamics and correction cells updated the latent state over the input sequence; a frame without any observed cells advanced the dynamics without correction. Decoder skip features came from the latest input frame containing observations. The decoder first estimated the current surface map; the latent state was then advanced six times without new observations to produce successive forecasts, with the skip features held fixed.

Training minimized a weighted sum of mean absolute errors over valid current targets, valid future targets, and target-valid cells hidden at the final input frame. The hidden-region term included both current and future predictions and received a weight of 2. The network used 16 base channels, a $3 \times 3$ recurrent kernel, a batch size of eight, and the Adam optimizer with a learning rate of 0.001 and gradient-norm clipping at 1.0. Training used only the 16 training recordings; the two validation recordings were used solely for checkpoint selection. A reference model was trained with standard masks for at most 30 epochs, with early stopping after six epochs without improvement in the validation objective. All other servo models were trained for 100 epochs without early stopping. In every case the checkpoint with the lowest validation objective was retained.

Five servo models were compared: the 30-epoch reference, two 100-epoch models trained with standard masks, and two 100-epoch models trained with mixed masks, where each pair contained one model trained from random initialization and one fine-tuned from cloth-pretrained weights. The mixed-mask model with the lower validation objective was designated the primary model. Each configuration was trained once with the same seed; because GPU arithmetic was not bit-reproducible, variation between repeated runs was not measured.

Cloth pretraining used the 12 training and 3 validation cloth sequences (10,295 training windows); the remaining three sequences were not used. The depth maps were aggregated onto a $32 \times 32$ camera-coordinate grid with 20~mm cells covering $640 \times 640$~mm, accepting depths of 400--2,000~mm and requiring at least three points per valid cell. Nominal timestamps were assigned from the frame indices at the published 5 frames per second. Windows contained five input frames (800~ms) and one 200~ms target, standard masks moved at every frame, and training ran for at most 50 epochs with early stopping after eight epochs without improvement. Fine-tuning on the servo recordings copied only the network weights: normalization was recomputed from the servo training data, the optimizer was reinitialized, and the learning rate was reduced to 0.0003.


\subsubsection{Synthetic Occlusion on Held-Out Motions}
Held-out evaluation used the unmasked recorded measurements as reference and reported mean absolute error and root-mean-square error separately for current reconstruction and each forecast step, with the sixth step corresponding to the nominal 200~ms horizon. Errors were also grouped by whether a cell was visible or unobserved at the input origin; the unobserved region included natural missing cells as well as synthetic masking. A last-observation baseline carried forward the most recent valid value within the input history. Because this baseline could not predict a cell never observed in the window, its errors were compared with model errors on the same supported cells; model errors were additionally reported over all unobserved cells with valid targets. The test windows overlapped in time and were not independent trials. This experiment evaluated synthetic missing observations and did not establish performance behind an actual nozzle, where tool depth returns, shadows, and other sensor artifacts may be present.

Standard masks were used for this test. To evaluate head-scale occlusion with true labels, the same 1,740 test windows were also scored with large stationary masks, constructed as above with masked fractions of 70--90\% of the grid and held fixed throughout the input history. Every hidden cell was then hidden throughout the window, so the last-observation baseline had no support. All five models were evaluated on identical windows and masks.

\subsubsection{Real Print-Head Occlusion}
The five reserved patterns were used to examine the model with a physical print head in view. Both recordings of each pattern were projected onto the same $23 \times 45$ grid with the same crop and depth gate as the training data, after verifying each decoded frame against its recorded checksum; the depth intrinsics and depth scale were confirmed to match those of the training recordings.

The print head was not segmented frame by frame. Instead, a conservative stationary mask was constructed from a hand-inspected bounding box of the head in the color image, padded by 20 pixels. For each grid cell, a depth envelope was taken from at most 120 evenly sampled frames of each training recording and widened by 20~mm; points spanning each cell and its envelope were projected into the color image using the factory depth-to-color extrinsics and color intrinsics, and cells whose projection overlapped the padded box were masked. Cells never observed in the sampled training frames were also masked. No measurements from the reference recordings were used to define the mask or the model inputs. With this construction, 825 of the 1,035 grid cells were masked, leaving approximately 15--18\% of all cells observed; the corresponding counts for 0-, 10-, and 30-pixel padding were 610, 720, and 930 cells.

Each blocked frame was paired with the reference frame nearest in elapsed time from the start of its own recording, within 20~ms. Using the blocked history alone, the model estimated the current surface and the surface six frames later. Windows whose input intervals fell outside the training tolerance were excluded, and forecasts were scored only when the camera timestamps through the future frame satisfied the same interval tolerance, and no future observations were supplied to the model. Differences from the paired reference were reported within the masked region and over visible cells, together with a last-observation baseline on the same supported cells. Because the mask was stationary, the baseline had no support within the masked region. As a measure of run-to-run agreement before any prediction, the two recordings were also compared over cells visible in both.

The differences were also expressed as vertical differences in the printer frame. A camera-to-printer rigid transformation was fitted after the recordings from four printed reference dots with known printer coordinates: each dot center was located in the color image, metric depth at that center was obtained through the device SDK, and a rotation and translation without scale were fitted by singular value decomposition. The fit residual was 0.29~mm RMS (maximum 0.36~mm), and the maximum leave-one-out error was 0.70~mm; no independent check point was measured. The transformation assumes that the camera pose did not change between the recordings and the registration, which was made two to three days later. Under it, the reference recordings mapped to a surface at printer Z of 149--154~mm, within about 2$^\circ$ of level, consistent with the reference dots at 149.7--150.4~mm. Each predicted cell was mapped to printer coordinates, and its printer Z was compared with the reference surface's Z at the same printer XY, interpolated linearly over the reference frame's valid cells (nearest valid cell outside their footprint). Because the camera axis was inclined about 61$^\circ$ from vertical, a camera-depth difference corresponds to roughly half its value vertically on a level surface.

The reference recordings were separate repeats rather than simultaneous measurements of the hidden surface. Fixed scene features were shifted by approximately 18 pixels between the two runs, and an exploratory three-dimensional registration of the fixed scene did not pass spatial hold-out checks, so the grids were compared without registration. Differences therefore combine prediction error, viewpoint change, imperfect motion repeatability, temporal alignment, sensor noise, and mask error, and were not interpreted as hidden-surface accuracy. The model did not segment the nozzle.

